\chapter{推广运营}\label{s:outreach}

在科技圈里，把大学和政府机构贬为行动迟缓的恐龙很时髦，
但依我的经验，它们并不比规模相当的公司更差。
你当地的学校董事会、图书馆、以及市议员的办公室，
或许能提供场地、资金、宣传、
和你可能还没接触过的其他群体的联系，
以及一大堆别的有用东西；
认识他们能帮你短期解决或避开问题，
长远也能带来回报。

\seclbl{市场营销}{s:outreach-marketing}

有学术或技术背景的人常常以为，
\gref{g:marketing}{市场营销}就是吹嘘和误导。
实际上，
它是从别人的视角看事情、
理解他们的需求和想要什么、
并解释你能怎样帮到他们——简而言之，
就是怎么教他们。
本章会看如何用前几章的想法，
让人理解并支持你正在做的事。

第一步是弄清你在向谁提供什么，
也就是究竟是什么带来了志愿者、
资金、
以及你维持下去所需的其他支持。
答案往往反直觉。
例如，
大多数科学家以为论文是他们的产品，
但实际是他们的拨款申请，
因为那才是带来拨款的东西~\cite{Kuch2011}。
论文是说服人们资助那些申请的广告，
就像如今专辑是说服人们买音乐家演唱会门票和 T 恤的东西。

假设你的群体为离开职场几年后重新就业的人，
提供周末编程工作坊。
如果工作坊参与者付的钱够覆盖你的成本，
那么他们就是你的客户，工作坊就是产品。
另一方面，
如果工作坊是免费、或学习者只交一点象征性的钱来降低放鸽子率，
那么你真正的产品可能是以下东西的某种混合：

\begin{itemize}

\item
  你的拨款申请；

\item
  资助你的公司想招聘的、你工作坊的往届学员；

\item
  市长向市议会做的年度报告里、关于你工作坊的半页总结，
  用来显示她如何支持本地科技行业；
  或者

\item
  志愿者从教学中得到的个人满足感。

\end{itemize}

和课程设计一样（\chapref{s:process}），
市场营销的头几步，
是为可能对你做的事感兴趣的人创建画像，\index{学习者画像}
并弄清你能满足他们的哪些需求。
总结后者的一种办法，是为不同画像写\grefdex{g:elevator-pitch}{电梯演讲}{电梯演讲}。
这些演讲有一个被广泛使用的模板：

\newpage
\begin{longtable}{ll}
  面向 & \emph{目标受众} \\
  他们 & \emph{对现有可选方案的不满} \\
  我们的 & \emph{类别} \\
  提供 & \emph{关键好处}。 \\
  不同于 & \emph{其他选择} \\
  我们的项目 & \emph{关键的与众不同之处。}
\end{longtable}

\noindent
接着周末工作坊的例子，
我们可以为参与者用这段话术：

\begin{quote}

  面向\emph{离开职场几年后重新就业}、
  且\emph{仍负有家庭责任}的人，
  我们的\emph{入门编程工作坊}
  提供\emph{带现场托儿的周末课程}。
  不同于\emph{在线课程}，
  我们的项目\emph{让人有机会遇见处于人生同一阶段的人}。

\end{quote}

\noindent
再为公司里可能资助这些工作坊的决策者用这段话术：

\begin{quote}

  面向\emph{想招聘入门级软件开发人员}、
  却\emph{很难找到来自多元背景的成熟候选人}的公司，
  我们的\emph{入门编程工作坊}
  提供\emph{潜在的新人}。
  不同于\emph{大学招聘会}，
  我们的项目\emph{把公司和五花八门的候选人连接起来}。

\end{quote}

如果你不知道不同的潜在相关方为什么会对你做的事感兴趣，
就去问他们。
如果你知道，
也还是去问：
答案会随时间变化，
你可能会发现以前忽略的东西。

有了这些话术之后，
它们就该驱动你放到网站和宣传材料上的内容，
帮人们尽快弄清
你和他之间有没有可聊的东西。
（不过你大概\emph{不该}逐字照抄它们：
科技圈很多人这套模板见得太多了，
再看到一遍就会眼神发直。）

在写这些话术时，
记住学编程有很多理由
（\secref{s:intro-exercises}）。
成就感、
对自身生活的掌控、
以及属于一个共同体，可能比金钱更能激励人
（\chapref{s:motivation}）。
他们可能因为朋友在做、所以也来当志愿者教书；
同样，
一家公司可能说他们资助面向经济弱势高中生的课程，
是因为想在未来有一批更大的潜在员工，
但那位 CEO 实际上可能只是因为这是对的事才这么做。

\seclbl{品牌与定位}{s:outreach-branding}

\gref{g:brand}{品牌}是某人听到一个产品名字时的第一反应；
如果那个反应是「那是啥？」，
那你（还）没有品牌。
品牌之所以重要，是因为
人们不会去帮一个自己不知道、或不在乎的东西。

今天关于品牌的讨论大多聚焦于
如何在线建立认知。
邮件列表、
博客、
和 Twitter 都给了你接触人的方式，
但随着虚假信息量的增加，
人们对每一次单独的打扰越来越不在意。
这让\gref{g:positioning}{定位}变得愈发重要。
它有时被叫做「差异化」，
是让你的供给区别于他人之处——你电梯演讲里「不同于」的那一段。
当你接触已经熟悉你所在领域的人时，
应该强调你的定位，
因为那才是会抓住他们注意力的东西。

还有其他能帮到你建立品牌的事。
一件是用道具，
比如某个学习者用家里找来的废料做的机器人~\cite{Schw2013}，
或另一个学习者为他父母的养老院做的网站。
另一件是做一段短片——几分钟以内——展示
学习者的背景和成就。
两者的目的都是讲故事：
尽管人们总是向你要数据，
但他们相信并记住的是故事。

\begin{aside}{奠基神话}
  一个人或群体能讲的最动人的故事之一，
  是他们为什么、以及怎样开始的。
  你教的，是不是你希望有人教过你、却没人教的东西？
  是不是有某一个特定的人，你想帮他，从此打开了闸门？
  如果你的网站上没有一段以「从前……」开头的话，
  考虑加上一段。
\end{aside}

一个关键步骤，是让你的组织在在线搜索里能被找到。\index{可发现性!组织}
\cite{DiSa2014b}发现，
家长为课外计算课用的搜索词，
实际上搜不到那些课，
许多其他群体也面临类似的挑战。
关于如何让东西能被搜到（也就是\gref{g:seo}{搜索引擎优化}，SEO），
流传着大量民间说法；
鉴于 Google 近乎垄断的权力和缺乏透明度，
其中大部分归根结底是设法比那些防止人们操纵排名的算法快一步。

除非你资金非常充足，
你能做的最好的事，就是定期搜一搜你自己和你的组织，
看看出来什么，
然后读读\hreffoot{https://moz.com/learn/seo/on-page-factors}{这些指南}，
尽可能改进你的网站。
记住\hreffoot{https://xkcd.com/773/}{这张 XKCD 漫画}：
人们不想了解你的组织架构图、或来一次你网站的虚拟导览——他们想要你的地址、
停车信息、
以及大致了解你教什么、
什么时候教、
它会怎样改变他们的生活。

\begin{aside}{不是每个人都活在线上}
  这些例子假设人们能上网、
  假设群体有钱、有材料、有空闲时间、和/或技术技能。
  很多人并不具备——事实上，
  那些服务经济弱势群体的组织，几乎肯定不具备。
  （正如 Rosario Robinson 所说：「免费，是给那些付得起免费的人用的。」）\index{Robinson, Rosario}
  在那些情境里，故事比课程大纲更重要，
  因为故事更容易转述。
  同样，
  如果你希望接触到的人上网不如你频繁，
  那么学校的公告板、
  当地图书馆、
  临时服务中心、
  和杂货店，可能是接触他们最有效的方式。
\end{aside}

\seclbl{冷启动电话的艺术}{s:outreach-cold-call}

建个网站、然后指望人们找到它很容易；
在没有任何事先介绍的情况下给人打电话、或去敲他们的门，
就难多了。
然而，和站起来讲课一样，
它是一门可以学会的手艺。
下面是说服人做事的十条简单规则：

\begin{description}

\item[1：别。]
  如果你非要说动某人做某事，
  那多半说明他其实并不想做。
  尊重这一点：
  从长远看，把某件特定的事放着不做，
  几乎总比用内疚、或任何只会招来怨恨的阴损心理把戏去逼人做要好。

\item[2：善良。]
  我不知道是不是真有这么一本书叫
  \emph{《忍者销售大师的秘密招数》}，
  如果有，
  它大概会告诉读者：为潜在客户做点事，
  能制造一种欠人情的感觉，
  进而提高成交几率。
  那也许管用，但只一次管用，而且很下作。
  另一方面，
  如果你真心善良、
  因为这是好人该做的而去帮别人，
  你也许反而会激励他们也做好人。

\item[3：诉诸更大的善。]
  如果你一开口就谈「这对他们有什么好处」，
  你就是在暗示他们把和你的互动
  当成一场可以讨价还价的商业价值交换。
  相反，
  先解释你想让他们帮忙的事会怎样让世界更好，
  而且\emph{要真心这么想}。
  如果你提议的事不会让世界更好，
  那就提议一件更好的。

\item[4：从小处开始。]
  大多数人理所当然地不愿一头扎进事情里，
  所以给他们一个试水的机会，
  让他们了解你和所有参与其中的人。
  如果事情就到此为止，别惊讶、也别失望：
  每个人都很忙、很累，有自己的项目，
  或者只是对合作该怎么进行有不一样的心智模型。
  记住 90-9-1 法则——90\% 的人会旁观，
  9\% 会开口，
  只有 1\% 会真的动手——并据此设定你的预期。

\item[5：不要建一个项目，要建一个共同体。]
  我以前属于一支其实从不打棒球的棒球队：
  我们的「比赛」只是大家聚在一起、享受彼此陪伴的借口。
  你大概不用做到那么极端，
  但和人一起喝杯茶、或庆祝人家第一个孙辈出生，
  能给你带来任何合理数量的钱都买不到的东西。

\item[6：建立一个连接点。]
  「我和 X 聊过」或「我们在 Y 见过」给了他们语境，
  这又让他们更自在。
  这必须具体：
  垃圾邮件发送者和冷启动电话已经把我们所有人训练成忽略任何以
  「我最近看到你的网站{\ldots}」开头的东西。

\item[7：把你要什么说具体。]
  人们需要知道这一点，
  才能判断他们拥有的时间和技能
  是否匹配你的需要。
  一开始就实事求是，也是尊重的表现：
  如果你说你只需要帮手搬几个箱子，
  实际上却是打包一整座房子，
  他们大概不会帮你第二次。

\item[8：建立你的可信度。]
  提一提你的支持者、
  你的规模、
  你的群体存在了多久，或者你过去取得的某项成就，
  好让他们相信你值得被认真对待。

\item[9：制造一点点紧迫感。]
  「我们希望在春天启动这个」比「我们最终想启动这个」更可能得到正面回应。
  不过，「一点点」这个词很重要：
  如果你的请求显得紧急，
  大多数人会以为你组织混乱、或出了什么岔子，
  然后可能就偏向谨慎。

\item[10：听懂言外之意。]
  如果你第一个求助的人说不，
  就去问别人。
  如果第五个或第十个人说不，
  就问自己：你想做的事到底讲不讲得通、值不值得做。

\end{description}

下面这封邮件模板遵循了所有这些规则。
它效果相当不错：
我们发现大约一半的邮件得到了回复，
其中大约一半想再多聊聊，
而其中大约一半带来了工作坊，
也就是说 10—15\% 的目标邮件转化成了工作坊。
如果你不习惯，这仍然可能相当打击人，
但比大多数组织对冷启动电话所预期的 2—3\% 回复率好得多。

\begin{quote}

  \noindent
  你好 NAME：

  希望你不介意我突然来信，
  但我想接着我们在 VENUE 的谈话往下谈，
  看看你有没有兴趣让我们办一期教师培训工作坊——我们
  正在安排接下来几周里的下一批。

  这个一天的工作坊会教你的志愿者们
  一些实用的、基于证据的教学实践。
  它以各种形式在六大洲为非营利组织、图书馆和公司办过一百多次，
  所有材料都可以在 http://teachtogether.tech 免费在线获取。
  主题会包括：

  \begin{itemize}
  \item 学习者画像
  \item 不同种类学习者之间的差别
  \item 用形成性评价诊断误解
  \item 教学作为表演艺术
  \item 什么激励、什么打击成人学习者
  \item 包容性的重要性，以及如何做一个好的盟友
  \end{itemize}
  
  如果这听起来有意思，
  请跟我说一声——我很乐意找机会聊聊方式方法。

  谢谢，

  NAME

\end{quote}

\begin{aside}{转介绍}
  和你工作相关的其他群体建立同盟，
  会在很多方面带来回报。
  其中之一是转介绍：
  如果某个来找你帮忙的人，由别的组织来服务会更好，
  那就花点时间做一次介绍。
  如果你做过好几次了，
  就在网站上加上点东西，帮下一个人找到他们需要的东西。
  你帮过的组织，很快也会反过来帮你。
\end{aside}

\seclbl{学术界的变革}{s:outreach-schools}

人人都害怕未知、害怕出丑。
结果就是，
大多数人宁愿失败也不愿改变。
例如，Lauren Herckis 研究了\index{Herckis, Lauren}
\hreffoot{https://www.insidehighered.com/news/2017/07/06/anthropologist-studies-why-professors-dont-adopt-innovative-teaching-methods}{为什么大学教师不采用更好的教学方法}。
她发现主要原因是害怕在学习者面前显得愚蠢；
次要原因包括
担心改变教学方法时不可避免的磕碰会影响课程评价
（这又可能影响晋升和终身教职）、
以及人们想继续效仿那些曾激励过自己的老师。
去争论这些问题是不是「真实」的没意义：
教职员工相信它们是真的，
所以任何与他们合作的计划都必须处理它们\footnote{
  还有教职员工在教学中普遍存在的固定型思维，
  也就是相信有些人「天生就是更好的老师」。}。

\cite{Bark2015}做了一项分两部分的研究，考察计算机科学教育者作为个人、组织、以及整个社会，
是如何采用新教学实践的。
他们提出并回答了三个关键问题：

\begin{description}

\item[教职员工是怎么听说新教学实践的？]
  他们主动寻找新做法，
  因为他们有解决问题的动力（尤其是学生参与度）、
  因为所在机构的有意推动而得知、
  从同事那里学来、
  或者从会议上预期\emph{和意外}的互动里得来
  （与教学有关或无关）。

\item[他们为什么去试？]
  有时是因为机构激励
  （例如他们创新以提高晋升机会），
  但在研究型大学里往往存在张力，
  因为关于教学重要性的说辞在很大程度上不被相信。
  另一个重要原因是他们自己的成本/收益分析：
  这项创新能帮他们省时间吗？
  第三个原因是他们受榜样激励——再说一遍，
  这主要影响那些旨在改善参与度和动机、
  而非学习成效的创新——第四个原因是信任的来源，
  比如他们在会议上遇到的、处在和自己相同处境、
  并报告了成功采用的同行。

  但教职员工的顾虑，常常没被倡导变革的人处理。
  第一个是格拉斯定律：
  任何新工具或新做法一开始都会拖慢你，
  所以新做法虽然长远看可能让教学更有效，
  短期内却负担不起。
  另一个是教室的物理布局让许多新做法很难实施：
  例如，
  讨论小组在阶梯教室里效果不好。

  但最能说明问题的结果是这一句：
  「尽管本身就是研究者，
  我们访谈的那些计算机科学教职员工，大多并不相信
  教育研究的结果是尝试新教学实践的、可信的理由。」
  这与其他发现一致：
  即便整个职业生涯都奉献给研究的人，也常常忽视教育研究。

\item[他们为什么继续用？]
  正如~\cite{Bark2015}所说：「学生的反馈至关重要，」
  而且往往是继续采用某项做法最强的理由，
  尽管我们知道学习者的自我报告与学习成效并不强相关~\cite{Star2014,Uttl2017}
  （不过讲座的到课率是参与度的好指标）。
  继续采用某项做法的另一个理由是机构要求，
  不过如果这是唯一的动机，
  一旦明确的激励或监督被撤掉，人们往往会放弃它。

\end{description}

好消息是，你可以系统地处理这些问题。
\cite{Baue2015}考察了美国退伍军人管理局内新医疗技术的采用。
他们发现，医学中基于证据的实践
平均要 17 年才能被纳入常规的全科实践，
而且这类实践里只有大约一半曾被广泛采用。
这一令人沮丧的发现促成了
\gref{g:implementation-science}{实施科学}的兴起，
它研究的是如何让人采纳更好的做法。

正如\chapref{s:community}所说，
起点是弄清你想帮的人认为自己需要什么。
例如，
\cite{Yada2016}总结了 K-12 教师对所要的准备和支持的反馈。
虽然它未必都适用于所有情境，
但和几个人喝杯茶、先听后说，
对他们愿不愿意尝试新东西有天壤之别。

一旦你知道人们需要什么，
下一步就是在机构自身的框架内逐步做出改变。
\cite{Nara2018}描述了一个基于紧密小群体和行政支持的三年制强化本科项目，
把毕业率提高到了三倍；
而~\cite{Hu2017}描述了为现有想教计算的高中教师
引入一个六个月认证项目的影响。
计算机教师的人数从 2007 年到 2013 年一直稳定，
但在引入新认证项目之后翻了两番，
而且没有稀释质量：
在教入门课程上，新转来教计算的教师，
似乎和受过更多计算训练的教师一样有效。

更广泛地说，
\cite{Borr2014}对高等教育中促成变革的方式做了分类。
这些类别由变革是个体性的还是系统性的、
以及它是指令性的（自上而下）还是涌现性的（自下而上）来界定。
试图促成（并让之持续）变革的人，在每种情境里扮演不同的角色，
也应采取不同的策略。
那篇论文接着详细解释了每一种方法，
而~\cite{Hend2015a,Hend2015b}以更可行动的形式呈现了同样的想法。

从外部来看，
你起初大概会落在「个体/涌现」这一类，
因为你会一个一个地接触老师，
试图自下而上地促成变革。
如果是这样，
Borrego 和 Henderson 推荐的策略，核心是
让老师单独或成组地反思自己的教学。
用现场编程向他们展示你做什么、或你用什么例子，
然后让他们也现场编程，
展示他们会如何在各自情境里运用那些想法和技巧，
这就给了每个人一个机会，去捡起对自己语境有用的东西。

\seclbl{散养式教学}{s:outreach-free-range}

学校和大学不是人们学编程的唯一去处；
过去几年，越来越多的人转向散养式工作坊
和强化集训营。
后者通常持续一到六个月，
由志愿者群体或营利性公司运营，
面向想转行进入科技行业的人。
有些质量很高，
但另一些存在的主要目的就是把钱从人身上拿走~\cite{McMi2017}。

\cite{Thay2017}访谈了这类集训营的 26 位往届学员——
这些集训营为那些早先错过计算机教育机会的人提供第二次机会
（不过这种说法，一碰到来自代表性不足群体的人，就带着相当大的假设了）。
集训营参与者面临巨大的个人成本和风险：
他们在入营前、中、后都必须投入大量时间、金钱和精力，
而转行可能得花一年甚至更久。
好几位受访者觉得，自己的证书被雇主看不起；
正如一些人所说，
拿到工作意味着通过面试，
但既然面试官往往不肯说出拒绝的理由，
就很难知道该改什么、或还该学什么。
许多人只好去实习（带薪或不带薪），
并花大量时间积累作品集和建立人脉。
他们最清楚地识别出的三个非正式障碍是行话、
冒充者综合征、
以及一种「融不进去」的感觉。

\cite{Burk2018}更深入地挖了挖这件事，
把科技行业招聘者想要的技能和资质，
与四年制学位和集训营所提供的东西做了比较。
基于对来自各种规模公司的 15 位招聘经理的访谈、以及一些焦点小组，
他们发现招聘者一致强调软技能
（尤其是团队合作、沟通、以及持续学习的能力）。
许多公司要求四年制学位
（不过未必是计算机科学的），
但许多公司也称赞集训营毕业生更年长、更成熟，
并有更新的知识。

如果你要去接触一个现有的集训营，
你最好的策略可能是强调你在教学上懂什么，
而不是你在技术上懂什么，
因为他们许多创办者和员工有编程背景，
却几乎没有或完全没有教育方面的训练。
本书前几章过去在这个受众里反响不错，
而~\cite{Lang2016}描述了
只需最小投入、低成本就能落地的基于证据的教学实践。
这些也许不是最有影响力的，
但早期赢下几局，有助于为更大的努力积累支持。

\seclbl{结语}{s:outreach-final}

你不可能独自改变大型机构：
你需要盟友，
而要有盟友，
你需要策略。
我发现的最有用的指南是~\cite{Mann2015}，
它编录了四五十种策略，
并按它们最适合在早期、后期、整个变革周期、还是遇到阻力时使用来组织。
其中一些模式包括：

\begin{description}

\item[放进他们的空间：]
  通过在整个组织里放置提醒，
  让新想法保持可见。

\item[代币：]
  为了让一个新想法留在一个人的记忆里，
  发放能让人联想到所引入主题的代币。

\item[支持型怀疑者：]
  请那些对你的新想法持怀疑态度的、有影响力的意见领袖，
  扮演「官方怀疑者」。
  用他们的评论来改进你的努力，
  即便你没能改变他们的想法。

\item[未来的承诺：]
  如果你能预见到自己的一部分需求，
  就可以向忙碌的人索取一份未来的承诺。
  如果给他们一些提前量，
  他们可能更愿意帮忙。
  
\end{description}

最重要的策略是
愿意根据你从想帮的人那里学到的东西，改变你的目标。
教他们怎么用电子表格的教程，
可能比 JavaScript 入门更快、更可靠地帮到他们。
我常犯的一个错误，是把让我自己充满热情的东西，
和别人应该知道的东西混为一谈；
如果你真想做一个伙伴，
永远记住：学习和改变都得是双向的。

建立关系最难的部分是起步。
每个月留出一两个小时，
去结识盟友、并维护你和他们的关系。
一个办法是请他们给建议：
他们觉得你该怎样提高人们对你所做之事的认知？
他们在哪里找到了上课的场地？
他们认为哪些需求没被满足、
而你能不能去满足？
任何存在了几年以上的群体都会有有用的建议；
被请教也会让他们感到受宠若惊，
而且下次你打电话时，他们会知道你是谁。

正如~\cite{Kuch2011}所说，
如果你没法在一个类别里当第一，
就试着创造一个新类别、让自己在其中当第一。
如果连这也不行，
就加入一个已有的群体，或者考虑干脆做点别的事。
这不是认输：
如果别人已经在做你心里想的事，
你要么出一份力，要么去做另一件你本可以做、而同样有用的事。

\seclbl{练习}{s:outreach-exercises}

\exercise{向市议员做推介}{个人}{10}

本章描述了一个
为重新就业的人提供周末编程工作坊的组织。
为这个组织写一段电梯演讲，
对象是该组织需要其支持的市议员。

\exercise{推介你的组织}{个人}{30}

找出你的组织需要其支持的两类人，
为每一类各写一段电梯演讲。

\exercise{邮件主题}{两人一组}{10}

为三封邮件写主题行（只写主题行）：
一封宣布一门新课程，
一封宣布一个新赞助方，
一封宣布项目领导层变动。
把你的主题行和同伴的比较，
看看能不能在合并各自最佳特点的同时把它们缩短。

\exercise{处理消极抵抗}{小组}{30}

不想改变的人有时会明说，
但往往也会用各种消极抵抗，
比如一次次就是拖着不做，
或者一个接一个地提出可能的问题，
让改变显得比实际更危险、更昂贵
\cite{Scot1987}。
以小组为单位，
列出三四个人们可能不想让你的教学计划推进的理由，
并说明在你现有的时间和资源下，你能做些什么来抵消每一个。

\exercise{为什么要学编程？}{个人}{15}

重温\secref{s:intro-exercises}里「为什么要学编程？」那道练习。
你教书的理由和学习者学习的理由，在哪里一致？
在哪里不一致？
这对你的市场营销有什么影响？

\exercise{对话型程序员}{思考—结对—分享}{15}

\gref{g:conversational-programmer}{对话型程序员}
是那些需要懂足够多的计算知识、
以便和程序员进行有意义交流、但自己并不打算编程的人。
\cite{Wang2018}发现，大多数学习资源都没照顾到这个群体的需求。
两人一组，
为一个半天工作坊写一段推介，对象是符合这一描述的人，
然后把你们这组的推介分享给全班。

\exercise{协作}{小组}{30}

先自己回答下面的问题，
再把答案和组里其他成员给的比较。

\begin{enumerate}

\item
  你和其他群体有任何协议或关系吗？

\item
  你想和其他任何群体建立关系吗？

\item
  有（或没有）协作会怎样帮你实现目标？

\item
  你的关键协作关系是什么？

\item
  这些是实现你目标的合适协作者吗？

\item
  你希望你的组织和哪些群体或实体建立协议或关系？

\end{enumerate}

\exercise{教育化}{全班}{10}

\cite{Laba2008}探究了为什么美国和其他国家
不断把社会问题的解决推给教育机构，
以及为什么这持续行不通。
正如他指出的：
「[教育]在促进种族、阶级和性别平等方面几乎没什么作为；
在提升公共健康、经济生产力和良好公民素养方面也一样；
在减少青少年性行为、交通事故死亡、肥胖和环境破坏方面也是。
事实上，
它在许多方面对这些问题产生了负面影响，
因为它把本可能产生更实质影响的资金和精力，从社会改革中抽走了。」
他接着写道：

\begin{quote}

  那么，面对这个机构没能做到我们对它的要求，
  我们该怎么理解它的成功？
  一种理解方式是：
  教育也许没在做我们要求的事，
  但它在做我们想要的事。
  我们想要一个机构，能以一种契合自由主义理想核心的个人主义的方式，
  去追求我们的社会目标，
  通过设法改变一个个学生的心、头脑和能力，来解决社会问题。
  换种说法，
  我们想要一个机构，通过它我们可以表达自己的社会目标，
  又不违反位于社会结构中心的个人选择原则，
  哪怕代价是达不到这些目标。
  所以教育可以充当公民自豪感的来源、
  我们理想的展示橱窗、
  以及一个就美好生活的不同愿景进行振奋人心、却终究无关紧要的争论的平台。
  同时，
  它也可以充当一个方便的替罪羊，
  让我们可以为它没能实现我们作为社会对自身的最高期望而责怪它。

\end{quote}

在学校里教计算思维和数字公民的努力，
如何契合这个框架？
集训营是避开了这些陷阱，还是只是换了个新面目把它们递送出来？

\exercise{机构层面的采用}{全班}{15}

重读\secref{s:outreach-schools}里给出的
采用新做法的动机清单。
其中哪些适用于你和你的同事？
哪些与你的情境无关？
当你和正规教育机构里的人打交道时，
你会强调哪些？

\exercise{如果一开始没成功}{小组}{15}

W.C.~Fields 大概从没说过：
「如果一开始没成功，就试了再试。
然后就放弃吧——一直当个傻瓜没好处。」
这仍是好建议：
如果你想接触的人没有回应，
那可能意味着你永远说服不了他们。
以 3—4 人为一组，
列一份简短的清单，写出哪些迹象表明你该停止去做某件你信仰的事。
其中有多少已经成真了？

\exercise{让它失败}{个人}{15}

\cite{Farm2006}给出了一些半开玩笑的规则，用来确保新工具\emph{不被}采用，
它们全都适用于新教学实践：

\begin{enumerate}

\item
  把它设为可选。

\item
  在培训上省钱。

\item
  别在真实项目里用它。

\item
  绝不把它整合进流程。

\item
  偶尔才用它。

\item
  把它作为某个质量计划的一部分。

\item
  把倡导者边缘化。

\item
  利用早期的失误。

\item
  只做小投入。

\item
  利用恐惧、不确定、怀疑、懒惰和惰性。

\end{enumerate}

其中哪些你最近见人做过？
哪些你自己做过？
它们以什么形式出现？

\exercise{导师指导}{全班}{15}

\hreffoot{http://www.iaamcs.org/}{非裔美国人计算机科学辅导研究所}
发布了\hreffoot{http://iaamcs.org/guidelines}{指导博士生的指南}。
各自读一遍，
然后作为一个班一起过一遍，
给你们自己群体的努力打分：+1（肯定在做）、
0（不确定或不适用）、
或 -1（肯定没在做）。
